% Generated draft: requires mathematical and source review.
\documentclass[11pt]{article}
\usepackage[margin=25mm]{geometry}
\usepackage{fontspec}
\setmainfont{Latin Modern Roman}
\usepackage{amsmath,amssymb,mathtools,unicode-math}
\setmathfont{Latin Modern Math}
\usepackage{enumitem,needspace,xurl,hyperref}
\setlength{\parindent}{0pt}
\setlength{\parskip}{5pt}
\setlength{\emergencystretch}{3em}
\newcommand{\N}{\mathbb{N}}
\newcommand{\Z}{\mathbb{Z}}
\newcommand{\Q}{\mathbb{Q}}
\newcommand{\R}{\mathbb{R}}
\newcommand{\C}{\mathbb{C}}
\newcommand{\F}{\mathbb{F}}
\newcommand{\A}{\mathbb{A}}
\newcommand{\PP}{\mathbb{P}}
\newcommand{\E}{\mathbb{E}}
\newcommand{\eps}{\varepsilon}
\newcommand{\dd}{\,\mathrm{d}}
\begin{document}
\section*{UnsolvedMath ID: 10400229}
\textbf{Source problem code:} catalogue entry
\subsection{Definitions and mathematical statement}
Let \(K\subset S^3\) be a prime knot, meaning a nontrivial knot not expressible as a connected sum of two nontrivial knots. Let \(\overline K\) be its mirror image. Say that \(K\) can be transformed into \(\overline K\) by one crossing change if some diagram of \(K\), after switching one crossing, represents \(\overline K\). This is an existence condition on a one-crossing-change presentation; it does not require the minimal crossing-change distance to be exactly \(1\). A knot is achiral if it is equivalent to its mirror under the intended knot-isotopy convention.

A. Stoimenow asks whether every such prime knot is achiral or ``(algebraically?) slice.'' Thus the printed conclusion contains an unresolved alternative. The two natural readings are
\[
 \bigl(K\text{ prime and }K\xrightarrow{\text{one crossing change}}\overline K\bigr)
 \ \Longrightarrow\ 
 \bigl(K\text{ achiral or }K\text{ slice}\bigr),
\]
or the same implication with \(K\) algebraically slice in place of slice. Here geometric sliceness means bounding an embedded disk in \(B^4\), in a smooth or locally flat topological category that must be specified. Algebraic sliceness means that a Seifert form is metabolic: for a genus-\(g\) Seifert surface there is a rank-\(g\) direct summand of its first integral homology on which the Seifert pairing vanishes.

The excerpt does not choose between these readings, specify the geometric slicing category, or clarify whether the word ``or'' is a disjunctive conclusion or presents alternative proposed conclusions. Those distinctions are retained rather than resolved by imposing the stronger statement.
\subsection{Short English statement}
If a prime knot becomes its mirror after changing one crossing, must it be achiral or slice, with algebraic sliceness explicitly offered as a possible alternative in the source?
\subsection{Sources}
\begin{itemize}
\item A. Stoimenow, attributed Problem 12.24 in T. Ohtsuki, \emph{Problems on invariants of knots and 3-manifolds} (2002), supplied PDF page 169. \url{https://msp.org/gtm/2002/04/gtm-2002-04-024s.pdf}. Recorded access 2026-07-29; extraction pdf-font-structure. No separate reference is supplied.
\item Local UnsolvedMath registry, \texttt{data/unsolvedmath/problems.json}, permanent ID \(10400229\), exact source code \texttt{catalogue entry}; research classification OPEN-TRIAGE.
\end{itemize}
\end{document}
